\chapter{Three Meanings of Mathematical Success}
\label{ch:success}

A report that a mathematical task has succeeded can mean at least three things, and the failure to distinguish them is the source of a large part of the confusion that this monograph addresses. The distinctions are not subtle once stated, and they have an exact form that the later chapters use. This chapter states them, proposes an initial taxonomy of the properties that a mathematical output may have, and explains why those properties cannot in general be reduced to one binary criterion of acceptance.

\section{Result, derivation, task}

\emph{Result success} is the property that the final output of a procedure is correct, meaning that it agrees with the right answer. A numerical computation that returns the right number has result success. \emph{Derivational success} is the property that the route by which the output was obtained is valid, meaning that each step follows from the steps before it by accepted inference. \emph{Task success} is the property that the output answers the question that was originally posed, in the sense in which it was posed.

The three come apart in each direction. A correct numerical answer may be reached through invalid reasoning: two errors may cancel, or a wrong method may land on the right value, as in the trajectory that computes $2 + 3$ as $2 \cdot 3 - 1$. A correct theorem may be proved through a derivation different from the one supplied, as in the symmetric-matrix case of Chapter~\ref{ch:proofnot}, so that derivational success of the produced proof does not carry derivational success of the supplied one. And a formally correct result may fail to answer the question originally posed, as in the substitution of $x^3 = 27$ for $x^2 = 9$ with $x > 0$, over the real numbers, where every step is valid and the answer $3$ is correct for both problems and yet the problem solved is not the problem posed.

The relation among the three is better understood as a system of predicates, none of which is in general a function of the others, than as a hierarchy. Let a trajectory $h$ produce an output $y(h)$ for a task $T$. Result success is $y(h) \in \mathrm{Sol}(T)$ when answers are checked against a solution set. Derivational success is the validity of $h$ as a derivation. Task success is the satisfaction of the obligations of $T$ by $h$, which includes that the derivation addresses $T$ and not a substitute. Result success follows from task success for the class of tasks whose obligations include correctness of the answer, and does not follow for tasks whose obligations omit it; the converse fails by the examples above, for every task class considered. Derivational success does not entail task success, since a valid derivation may be of the wrong problem, and task success does not entail derivational success of every step if the task is satisfied by a different route.

\section{A taxonomy of properties}

The three notions are the coarsest of a family. A usable taxonomy distinguishes at least the following properties of a mathematical output relative to a task. \emph{Correctness} is the truth of the stated result. \emph{Validity} is the correctness of the inferences of the derivation from its stated premises, which does not include the truth of those premises. \emph{Fidelity} is the correspondence of the stated or formalized result, and of the argument, to the intended statement and argument. \emph{Relevance} is the bearing of the output on the question asked, as distinct from its correctness. \emph{Completeness} is the coverage of all cases or conditions that the task requires. \emph{Adequacy} is the sufficiency of the output for the purposes for which it was requested, which may include being intelligible, being checkable, or being usable in later work.

The claim made about these properties is that none is in general a consequence of another, and it is made by examples and not by an exhaustive matrix of the thirty ordered pairs. The examples establish that the following properties can be had without the one that follows. A correct but irrelevant result is a true theorem offered in response to a different question, and has correctness without relevance. A valid but incomplete proof treats the cases it treats correctly and omits others, and has validity without completeness. A valid derivation from a false premise has validity without correctness. A faithful but incorrect formalization of a false statement is an accurate translation of an error, and has fidelity without correctness. An adequate but not fully valid argument is a sketch sufficient to convince an expert and containing a gap, and has adequacy without validity. The remaining ordered pairs are not claimed here, and some may be excluded in particular classes of task; the independence asserted is relative to the predicates as defined and to the existence of such instances. The taxonomy does not claim to be minimal or exhaustive, and its purpose is to show that the single word \emph{correct} in ordinary use stands for a conjunction of properties that can be separately satisfied.

\section{Why no binary criterion suffices}

A binary acceptance criterion is a function from outputs to $\{0,1\}$. If it is to be used as a proxy for the conjunction of the properties above, it must factor through them in the sense that the value depends only on their joint values and agrees with the intended aggregate. Two obstacles arise. The first is that the conjunction of several independent properties does not determine the value of any single one, so an acceptance criterion that is true exactly when all hold will be false in cases where only the property of interest to a particular user fails, and will conflate failures that call for different remedies. The second is the one made precise in Chapter~\ref{ch:reward} for rewards (RV-010): a reward that takes equal values on two outputs that differ in the task predicate cannot agree with the task on both, and the irreducible disagreement is measured by the conditional impurity within the level sets of the reward.

The consequence for reporting is that a verification should state which of the properties it establishes. The kernel establishes validity of the formal derivation relative to a stated environment, and it establishes correctness of the formal statement in the sense that the statement is provable there. It establishes nothing about fidelity, relevance, or adequacy. A report that says the proof was verified is correct if read as a statement about the first and misleading if read as a statement about the rest, and the monograph's practice is to attach to the word the property it names.

\section{Two results previewed}

Two of the monograph's results are stated informally here because they express the distinctions of this chapter and are proved later. The first is RV-004, which says that a proof can be valid under strengthened hypotheses without proving the original statement. If a derivation proves $H' \to G$ where $H'$ is strictly stronger than the intended hypothesis $H$, the derivation is valid and relevant to a special case, and it does not by itself yield a proof of $H \to G$, which is the stronger theorem. The second is the substitution result RV-015, proved in Chapter~\ref{ch:substitution}, which states, for a sound and complete certifier, exactly when acceptance for a substituted task transfers to the original. The proof of RV-004 is given in Chapter~\ref{ch:nondischarge}. Both are instances of the general fact that validity of a derivation is relative to what the derivation is a derivation of.
